\section{A whole-space bound for the good auxiliary labels}
\label{sec:good_auxiliary_label_entropy}

The logarithmic dimension comparison under restriction, together with the
actual occurring-label dimension bound on the bad copies, yields an
operator inequality before measuring a physical excitation component.

\begin{theorem}[Good auxiliary label bound]
\label{thm:good_auxiliary_label_entropy}
\lean{TensorPower.copyPerm_groupedCopies_labelEntropy_bounds}
\leanok
Let $C$ be a finite coordinate set, put $d=|C|$, and specify a split of the
$k$ copies into groups of sizes $m$ and $r$, where $m+r=k$. For the actual
copy-permutation action on $\mathbb C^{C^k}$, let $F_g,F_b,F_w$ denote the
central logarithmic label observables for the good, bad and whole groups.
Using the real logarithm extended by $\log0=0$, one has, on the whole space,
\begin{align}
 F_b&\le r\log d\,I,\\
 F_w&\le F_g+\left(r\log d+\log\binom{k}{r}\right)I.
 \label{eq:good_auxiliary_label_entropy}
\end{align}
Empty coordinate sets and zero copy groups are included.
\end{theorem}
\begin{proof}
\leanok
Every nonzero bad-label projection has irreducible dimension
$d_\beta\le d^r$. Positivity of $d_\beta$ permits taking logarithms,
which gives $\log d_\beta\le r\log d$. Summing over the actual orthogonal
resolution gives the first inequality; zero projections contribute nothing.
The subgroup operator bound gives
$F_w\le F_g+F_b+\log\binom{k}{r}\,I$. Adding the first inequality proves
the second one. The empty cases use the stated total logarithm convention.
\end{proof}

This is the operator input to
\cite[Section~7, equation \textup{(comparator:good-auxiliary)}]{OpenAI2026AreaLaw}.
To apply it to an excitation component, one still uses the actual
whole-label threshold and preservation of whole-copy auxiliary labels
under the physical excitation projection. No commutation of that projection
with the band metric is claimed here.
